% !TEX program = pdflatex
% ---------------------------------------------------------------------------
% likelihood_engine.tex -- how the likelihood engine works: the model it scores,
% every equation it evaluates, and how it is checked.
%
% Build: pdflatex likelihood_engine.tex  (twice, for cross-references).
% Code: analyses/2026-10_likelihood-engine/src/ (01 reference twin, 02 JAX engine,
% 03 validation, 04 timing). Design record: notes/sciphy_notes.md, H.8.
%
% Notation follows the house register (sciphy_notes.md, H.6.0): lambda is the
% EDITING rate and f is reserved for functions. A branch is named by its lower
% node (w), never b -- b is the division rate.
% ---------------------------------------------------------------------------
\documentclass[11pt]{article}
\usepackage[margin=1in]{geometry}
\usepackage{amsmath,amssymb,mathtools}
\usepackage{booktabs}
\usepackage{array}
\usepackage[hidelinks]{hyperref}

\newcommand{\E}{\mathbb{E}}
\newcommand{\Prob}{\mathbb{P}}
\newcommand{\Sil}{\mathsf{S}}               % the silenced state / slot
\newcommand{\LamT}{\Lambda_T}
\newcommand{\lcp}{\mathrm{lcp}}
\newcommand{\pa}{\mathrm{pa}}
\newcommand{\ch}{\mathrm{ch}}
\newcommand{\Obs}{\mathrm{Obs}}
\newcommand{\Out}{\mathrm{Out}}
\newcommand{\miss}{\textup{?}}
\newcommand{\Hbar}{\bar H}
\newcommand{\GamT}{\Gamma_T}
\newcommand{\sig}{\sigma}
\newcommand{\note}[1]{\par\smallskip\noindent\emph{#1}\par\smallskip}

\title{The likelihood engine:\\ an exact per-clone likelihood for DNA Typewriter lineage trees}
\author{Justin}  % edit
\date{Draft, \today}

\begin{document}
\maketitle

\noindent This document explains the likelihood engine in
\texttt{analyses/2026-10\_likelihood-engine/}: the model it scores, every quantity it computes,
why each equation has the form it does, and how the code is checked. It is the core of our own tree
reconstruction tool. The engine answers one question:

\begin{quote}
\emph{Given one clone's tree --- who descended from whom, and when each division happened --- and
the parameters of editing, silencing and readout, how probable are the tapes we actually read?}
\end{quote}

\noindent Everything else (finding good trees, putting error bars on them, inferring signalling
history) is built on top of that one function. Section~\ref{sec:prov} lists what is borrowed from
SciPhy, LAML, LAML-Pro and ConvexML, and what is new here.

% ===========================================================================
\section*{Notation}
% ===========================================================================
Symbols are grouped by what they describe. Most are re-defined where they are first used.

\begin{center}\small
\begin{tabular}{@{}lp{0.60\textwidth}l@{}}
\toprule
symbol & meaning & units / range \\
\midrule
\multicolumn{3}{@{}l}{\emph{Tree and time}}\\
$t$ & calendar time as a fraction of the experiment: clone founder at $t=0$, harvest at $t=1$ & $[0,1]$ \\
$h=1-t$ & height: time remaining until harvest & $[0,1]$ \\
$n$ & number of cells (tips) in the clone; the tree is rooted and binary with $n-1$ internal nodes & count \\
$o$ & the origin: the founder cell, at $t=0$ & --- \\
$r$ & the root: most recent common ancestor of all sampled cells, at $t_r\in(0,1)$ & --- \\
stem & the branch from $o$ to $r$, over $[0,t_r]$ & --- \\
$v,u,w$ & nodes; $\pa(v)$ is the parent, $\ch(v)$ the children & --- \\
branch $w$ & the branch \emph{into} node $w$, from $\pa(w)$ down to $w$, over $[t_{\pa(w)},t_w]$ & --- \\
$t_v$, $h_v$ & time and height of node $v$; every tip has $t=1$, $h=0$ & $[0,1]$ \\
$x_v$ & unconstrained coordinate of internal node $v$ (Section~\ref{sec:coords}) & real \\
\midrule
\multicolumn{3}{@{}l}{\emph{Tapes, cells, reads, states}}\\
$k$, $z$ & tapes per cell; $z=1,\dots,k$ indexes them & count \\
$N$ & slots per tape & $6$ \\
$M$ & number of distinct insert symbols; $i$ indexes them & count \\
$c$ & a cell (a tip of the tree) & --- \\
$d_{cz}$ & the read of tape $z$ in cell $c$: an observed string $v_{cz}=(s_1,\dots,s_m)$, $1\le m\le N$, or $\miss$ (missing) & --- \\
$X_{vz}$ & the true, hidden state of tape $z$ at node $v$: a string of at most $N$ symbols, or $\Sil$ (silenced) & --- \\
$d$ & a slot: the depth of a tape, i.e. how many symbols are written, $0\le d\le N$; plus the slot $\Sil$ & --- \\
\midrule
\multicolumn{3}{@{}l}{\emph{Editing}}\\
$\LamT$ & expected number of edits on an average tape over the whole experiment (if it never filled) & edits / tape \\
$\lambda_0(t)$ & time shape of the editing rate, $\ge0$, with $\int_0^1\lambda_0=1$ & --- \\
$W(t)=\int_0^t\lambda_0$ & share of the experiment's editing spent by time $t$ & $[0,1]$ \\
$r_z$ & speed of tape $z$ relative to the average tape & mean $1$ \\
$\lambda_z(t)=r_z\LamT\lambda_0(t)$ & editing rate of tape $z$ at time $t$ & edits / tape / experiment \\
$m_w$ & per-branch rate multiplier: a hook for lineage-varying editing; fixed at $1$ & --- \\
\midrule
\multicolumn{3}{@{}l}{\emph{Composition}}\\
$\xi_i(t)$ & probability that an edit made at time $t$ writes symbol $i$; $\sum_i\xi_i(t)=1$ & probability \\
$\xi^{L}$ & the background (lineage) composition, measured and fixed & probabilities \\
\midrule
\multicolumn{3}{@{}l}{\emph{Silencing}}\\
$\Hbar$ & expected fraction of cell $\times$ tape entries silenced by harvest & $[0,1)$ \\
$\GamT=-\ln(1-\Hbar)$ & total silencing exposure of one root-to-tip lineage & events / tape \\
$\zeta_0(t)$, $W_\zeta(t)$ & time shape of silencing ($\int_0^1\zeta_0=1$) and its cumulative & --- \\
$\Gamma(t)=\GamT W_\zeta(t)$ & silencing exposure accumulated by time $t$ & events / tape \\
\midrule
\multicolumn{3}{@{}l}{\emph{Readout}}\\
$\pi_{cz}=\sig(\alpha_c+\beta_z)$ & probability that tape $z$ of cell $c$ is lost technically; $\sig(y)=1/(1+e^{-y})$ & probability \\
$\alpha_c$, $\beta_z$ & cell and tape offsets (larger = worse captured / worse read) & logits \\
\midrule
\multicolumn{3}{@{}l}{\emph{Tree prior}}\\
$b$, $\delta$, $\rho$ & division rate, death rate (per experiment), fraction of the final population sampled & --- \\
\midrule
\multicolumn{3}{@{}l}{\emph{Computed by the engine}}\\
$\mu_w$, $G_w$ & expected edits on a tape, and silencing exposure, along branch $w$ & edits / tape; events / tape \\
$\lcp_{vz}$ & longest common prefix of the tapes read below node $v$ & string \\
$K_{wz}[d',d]$ & branch kernel: probability of going from slot $d'$ to slot $d$ along branch $w$ & probability \\
$L_{vz}(d)$ & partial likelihood: probability of everything read below $v$, if $v$ is in slot $d$ & probability \\
$\ell_z$ & log-likelihood of tape $z$ (the engine's output, one per tape) & nats \\
\bottomrule
\end{tabular}
\end{center}

% ===========================================================================
\section{The model the engine scores}\label{sec:model}
% ===========================================================================

The engine scores exactly the model our simulator generates (``rung 2'' of the simulator: editing,
heritable silencing, technical dropout and depth-0 censoring). It has four parts.

\paragraph{1. Tree and time.} One founder cell at $t=0$ divides at rate $b$ and dies at rate
$\delta$. At harvest, $t=1$, a fraction $\rho$ of the population is sampled. The tree relating the
sampled cells is rooted and binary, and every tip sits at $t=1$ (it is \emph{ultrametric}). Above the
root runs a \emph{stem} back to the founder. Time is measured in units of the experiment, so
calendar days are $t\times T$ with $T$ known.

\paragraph{2. Editing.} Every tape starts \emph{blank} at the founder. Along every lineage, and
independently for each tape, edits arrive at random --- a Poisson process with rate
$\lambda_z(t)=r_z\LamT\lambda_0(t)$. Each edit \emph{appends} one symbol to the tape, the symbol being
$i$ with probability $\xi_i(t)$. Once all $N$ slots are full, nothing more is recorded. At a
division both daughters inherit a copy of the tape.

\paragraph{3. Silencing.} Along every lineage, independently for each tape, the tape can be
silenced, at rate $\zeta(t)=\GamT\zeta_0(t)$. Once silenced, it stays silenced in every descendant and
is never read again (a \emph{Dollo} character: lost once, never regained).

\paragraph{4. Readout.} At harvest each cell's tapes are read.
\begin{itemize}\itemsep0pt
\item A silenced tape is always reported missing.
\item Otherwise the tape is lost technically with probability $\pi_{cz}=\sig(\alpha_c+\beta_z)$.
\item A tape that was captured but is still \emph{blank} is also reported missing. This is how Park's
  pipeline codes it (a tape whose first slot is empty is called \texttt{ABSENT}); we call it
  \emph{censoring}.
\item Otherwise the tape is read exactly.
\end{itemize}

\paragraph{Assumptions.} Given the tree and the parameters: tapes are independent of one another;
silencing is independent of what is written on the tape; technical losses are independent across
cells and tapes given $\alpha$ and $\beta$. The first is the strongest. A cell's tapes share its
chromatin state, its prime-editor supply and its cell cycle, so in reality they are not quite
independent; treating them as independent overstates how much evidence the data carry (rows A5/A8 of
the project's assumption table).

% ===========================================================================
\section{What a likelihood is, and why the engine works tape by tape}\label{sec:lik}
% ===========================================================================

\subsection{One story, and its probability}

Fix one tape $z$. At every node of the tree the tape had some true state --- a string of symbols, or
$\Sil$. We see the true states only at the tips, and only through the readout; the internal states
are hidden. Pick a complete \emph{story}: a true state for every node. Its probability is a product
of local pieces:
\begin{equation}
\underbrace{P_{\rm stem}(\varnothing\to X_{rz})}_{\text{founder blank; stem}}\;
\prod_{\text{branches }w}\underbrace{P_w\big(X_{\pa(w)z}\to X_{wz}\big)}_{\text{what happens along branch }w}\;
\prod_{\text{cells }c}\underbrace{\mathcal O_{cz}\big(d_{cz}\mid X_{cz}\big)}_{\text{what the sequencer reported}} .
\label{eq:story}
\end{equation}
Each factor answers one local question. A \emph{branch factor} asks: given the tape at the top of
this branch, how likely is its state at the bottom? An \emph{emission factor} $\mathcal O$ asks:
given a cell's true tape, how likely is what was reported?

\subsection{The likelihood is the sum over every story}

We do not know the hidden states, so the probability of the reads is the sum of~\eqref{eq:story}
over every possible story:
\begin{equation}
L_z=\sum_{\text{all hidden states}}\ \eqref{eq:story},\qquad \ell_z=\log L_z .
\label{eq:Lz}
\end{equation}
Because tapes are independent given the tree and parameters, the probability of all reads is the
product over tapes, and its logarithm is a sum:
\begin{equation}
L=\prod_{z=1}^{k}L_z,\qquad \ell=\log L=\sum_{z=1}^{k}\ell_z .
\label{eq:product}
\end{equation}
The engine returns the vector $(\ell_1,\dots,\ell_k)$, in nats; the clone's log-likelihood is its
sum. Listing stories is impossible --- with about 100 symbols and $N=6$ there are some $10^{12}$
possible states per node --- so Sections~\ref{sec:prune}--\ref{sec:kernel} compute the sum without
listing anything.

\subsection{Two ways of reading the same number}\label{sec:normalise}

The quantity $P(\text{reads}\mid\text{tree},\theta)$, where $\theta$ stands for all the parameters,
can be read in two ways, and only one of them adds up to 1.
\begin{itemize}
\item \textbf{As a function of the parameters} (reads fixed, $\theta$ varying) it is the
  \emph{likelihood}. It does \emph{not} integrate to 1 over $\theta$ or over trees, and nothing
  requires it to. This is the familiar statement that a likelihood is not a probability
  distribution.
\item \textbf{As a function of the reads} (tree and $\theta$ fixed, reads varying) it is the
  probability distribution of what an experiment could produce. Summed over \emph{every possible set
  of reads} --- every cell's tape reported as some string of length $1$ to $N$, or missing --- it
  must equal exactly 1, because some set of reads always occurs.
\end{itemize}
The second property is what validation check V2 tests. It is not a property we impose; it is
automatically true of a correct model, and it fails if, for example, a missing tape is counted under
two causes at once. That makes it a sharp test of the readout and censoring terms.

% ===========================================================================
\section{Pruning: computing the sum by passing messages up the tree}\label{sec:prune}
% ===========================================================================

Look at where one internal node's state, say $X_{uz}$, appears in~\eqref{eq:story}: only in the
branch into $u$ and the two branches out of it. Nothing else depends on it. So the sum over $X_{uz}$
can be done locally, before anything is decided higher up. Doing this at every node, from the tips
upward, is Felsenstein's pruning algorithm.

\paragraph{Partial likelihood.} At every node $v$ keep a vector with one entry per possible state:
\begin{equation}
L_{vz}(x)=\Prob\big(\text{everything read on tape }z\text{ below }v\ \big|\ X_{vz}=x\big).
\end{equation}
Think of it as $v$'s report card: ``if I had been in state $x$, this is how well that explains what
my descendants show.''

\paragraph{Messages.} A child $w$ sends its parent a message:
\begin{equation}
M_{w\to\pa(w),z}(x')=\sum_{x}P_w(x'\to x)\,L_{wz}(x).
\label{eq:message}
\end{equation}
In words: if the parent was in state $x'$, the child could have ended in any state $x$; weight each by
how likely the branch was to carry $x'$ to $x$, and by how well $x$ explains the child's own subtree.

\paragraph{Combining.} Once the parent's state is fixed, its two daughter lineages evolve
independently, so their messages multiply:
\begin{equation}
L_{vz}(x')=\prod_{w\in\ch(v)}M_{w\to v,z}(x').
\label{eq:combine}
\end{equation}

\paragraph{The root.} The founder is blank, so the last message comes up the stem from the blank
state:
\begin{equation}
L_z=\sum_{x}P_{\rm stem}(\varnothing\to x)\,L_{rz}(x).
\label{eq:root}
\end{equation}
One pass from the tips to the root computes~\eqref{eq:Lz} exactly. What remains is to say which
states each vector needs (Section~\ref{sec:slots}), what $P_w$ is (Section~\ref{sec:kernel}) and what
the tips start with (Section~\ref{sec:tips}).

% ===========================================================================
\section{Which states each node needs: depth slots}\label{sec:slots}
% ===========================================================================

Picture the space of all possible tape states as a tree of strings (a \emph{trie}): the blank tape at
the top, branching into every one-symbol string, each of those into every two-symbol extension, and
so on down to $N$ symbols. Editing only ever moves a tape \emph{down} this trie, one step per edit.
Three facts follow, and together they shrink every node's vector from $10^{12}$ entries to $N+2=8$.

\paragraph{Fact 1: a read tape pins its ancestors to one path.} If cell $c$'s tape was read as
$v_{cz}$, every ancestor of $c$ held a \emph{prefix} of $v_{cz}$ (editing only appends). If several
tapes were read below node $v$, $v$ held a prefix of all of them, i.e.\ a prefix of their
\emph{longest common prefix}
\begin{equation}
\lcp_{vz}=\text{longest common prefix of }\{v_{cz}:c\text{ below }v,\ d_{cz}\ne\miss\}.
\end{equation}
There are at most $N+1$ such prefixes, from the blank tape up to $\lcp_{vz}$ itself.

\paragraph{Fact 2: Dollo.} A silenced tape is never read and silencing is inherited, so a node with
any read tape below it cannot have been silenced.

\paragraph{Fact 3: a subtree with nothing read only constrains depth.} If every tape below $v$ is
missing, the reads say nothing about \emph{which} symbols $v$'s tape held --- only whether it was
blank (a blank tape can be censored) or silenced. Formally, $L_{vz}(x)$ depends on $x$ only through
its depth $|x|$ (and whether $x=\Sil$). The reason: a missing tip's emission depends only on
``blank or not''; how many edits a branch adds depends only on how many slots are open; and summing
over which symbols were written gives 1 for every write. This was checked against exhaustive
enumeration of all states on small trees (agreement to $3\times10^{-15}$).

\paragraph{Depth slots.} So every node, for every tape, needs a vector of $N+2$ numbers, indexed by
slot $d=0,1,\dots,N$ and $\Sil$, read in one of two ways:
\begin{itemize}
\item \textbf{anchored node} (some tape read below it): slot $d$ means the specific string
  $\lcp_{vz}[1..d]$, for $0\le d\le|\lcp_{vz}|$; slots beyond are impossible (zero); $\Sil$ is
  impossible;
\item \textbf{floating node} (nothing read below it): slot $d$ means ``any string of depth $d$'';
  $\Sil$ is allowed.
\end{itemize}
If a child is anchored, so is its parent, and the parent's lcp is a prefix of the child's --- so the
parent's slot $d'$ and the child's slot $d$ describe the same string up to depth $d'$.

\paragraph{Example.} A clone of four cells, $((c_1,c_2),(c_3,c_4))$, internal nodes $u$ (above
$c_1,c_2$), $w$ (above $c_3,c_4$) and the root $r$. On one tape, $c_1$ read $(A,B)$, $c_2$ read
$(A,C,A)$, and $c_3,c_4$ are missing. Then $u$ is anchored with $\lcp=(A)$ --- its tape was blank or
held $A$; $w$ is floating; the root is anchored with $\lcp=(A)$. The two possible states of $u$ are
two stories about the shared $A$: written before $u$ divided and inherited by both daughters, or
written independently on both daughters' branches (a coincidence, homoplasy, weighted by
$\xi_A$ twice). Pruning weighs both by their probability.

% ===========================================================================
\section{Time on a branch}\label{sec:time}
% ===========================================================================

A branch $w$ covers the calendar interval $[t_{\pa(w)},t_w]$. Two numbers summarise what can happen
to a tape on it.

\paragraph{Expected edits.} Integrating the editing rate over the branch,
\begin{equation}
\mu_{wz}=\int_{t_{\pa(w)}}^{t_w}\lambda_z(s)\,m_w\,ds=r_z\,\LamT\,m_w\,\big[W(t_w)-W(t_{\pa(w)})\big]
\quad\text{edits per tape.}
\label{eq:mu}
\end{equation}
With a constant rate, $\lambda_0\equiv1$ and $W(t)=t$, this is just
$r_z\LamT\times(\text{branch duration})$. With a time-varying rate, two branches of equal duration can
carry different amounts of editing, depending on where they sit in time --- which is why the engine
works with node \emph{times}, not just branch lengths.

\paragraph{Silencing exposure.} In the same way,
\begin{equation}
G_w=\Gamma(t_w)-\Gamma(t_{\pa(w)})=\GamT\,\big[W_\zeta(t_w)-W_\zeta(t_{\pa(w)})\big]
\quad\text{events per tape,}
\label{eq:G}
\end{equation}
and the probability that the tape is \emph{not} silenced along the branch is $e^{-G_w}$.

\paragraph{Only totals matter when the composition is fixed.} If the editing rate speeds up and slows
down within a branch but the symbol probabilities $\xi$ do not change, only the total $\mu_{wz}$
matters, not when within the branch the editing happened: speeding a clock up or down plays the same
film faster. The same holds for silencing and $G_w$. Breakpoints within a branch (\emph{segments})
are needed only where the composition $\xi$ changes, because then which symbol sits in a slot depends
on whether it was written before or after the change. In the first build the composition is
constant, so every branch is a single segment with the totals~\eqref{eq:mu}--\eqref{eq:G}.

% ===========================================================================
\section{The branch kernel}\label{sec:kernel}
% ===========================================================================

The kernel $K_{wz}[d',d]$ is the branch factor $P_w$ written on depth slots: the probability that tape
$z$ ends branch $w$ in slot $d$ given that it started in slot $d'$. For an anchored child it also
includes the probability of having written exactly the symbols $\lcp_{wz}[d'+1..d]$ on the way.

\subsection{One segment: three factors}

On a branch (or segment) with expected edits $\mu$, silencing exposure $G$ and composition $\xi$:
\begin{equation}
\boxed{\;K[d',d]=\underbrace{e^{-G}}_{\text{not silenced}}\times
\underbrace{C_N(d'\to d;\mu)}_{\text{right number of edits}}\times
\underbrace{\Xi(d',d)}_{\text{right symbols}},\qquad d'\le d.\;}
\label{eq:kernel}
\end{equation}

\paragraph{Right number of edits.} Edits arrive as a Poisson process, so the number attempted on
the branch is Poisson with mean $\mu$. If the tape does not fill ($d<N$), exactly $d-d'$ edits
happened:
\begin{equation}
C_N(d'\to d;\mu)=e^{-\mu}\,\frac{\mu^{\,d-d'}}{(d-d')!},\qquad d<N.
\end{equation}
If the tape fills ($d=N$), at least $N-d'$ edits were attempted --- any extra had nowhere to go:
\begin{equation}
C_N(d'\to N;\mu)=1-\sum_{j=0}^{N-d'-1}e^{-\mu}\frac{\mu^j}{j!},\qquad C_N(N\to N;\mu)=1 .
\end{equation}

\paragraph{Right symbols.} For an anchored child, each written slot must hold the symbol the lcp
says, each chosen independently:
\begin{equation}
\Xi(d',d)=\prod_{j=d'+1}^{d}\xi_{\lcp_{wz}[j]} .
\end{equation}
For a floating child, $\Xi=1$: whatever was written is summed over, and every write's probabilities
add up to 1.

\paragraph{Not silenced.} Silencing ignores what is on the tape, so it contributes the separate
factor $e^{-G}$. For a floating child the silenced slot gets the rest:
$K[d',\Sil]=1-e^{-G}$, and $K[\Sil,\Sil]=1$. For an anchored child the silenced column is never used
(Fact 2). Editing never removes symbols, so $K[d',d]=0$ when $d<d'$.

\subsection{Several segments}

If the composition changes along the branch, cut it at the changes into segments $k=1,\dots,K$, each
with its own $\mu_k$, $G_k$ and $\xi^{(k)}$. The tape passes the boundaries between segments at some
depths $j_1\le j_2\le\dots\le j_{K-1}$; write $j_0=d'$ and $j_K=d$. Then
\begin{equation}
K_{wz}[d',d]=e^{-\sum_kG_k}\sum_{d'\le j_1\le\cdots\le j_{K-1}\le d}\ \prod_{k=1}^{K}
C_N\big(j_{k-1}\to j_k;\mu_k\big)\,\Xi^{(k)}(j_{k-1},j_k),
\qquad K_{wz}[d',\Sil]=1-e^{-\sum_kG_k}.
\label{eq:multiseg}
\end{equation}
The sum runs over every way the tape's growth could be split among the segments; each split scores
segment $k$'s edits with segment $k$'s composition. Equivalently, $K_{wz}$ is the matrix product of
the segments' kernels in time order. When the composition is the same in every segment, the sum
collapses to one segment with $\mu=\sum_k\mu_k$ and $G=\sum_kG_k$, because Poisson counts add
(checked numerically, to $2\times10^{-16}$).

\paragraph{Segment boundaries are not inferred.} They are the knots of $\lambda_0$, $\zeta_0$ and of
the composition schedule --- fixed by the analyst or by the experimental design. What is inferred are
node times, and the kernel changes smoothly with them:
$\partial\mu_{wz}/\partial t_w=r_z\LamT m_w\lambda_0(t_w)$.

\paragraph{A cross-check.} The same kernel can be computed by writing the processes as rates on the
slots and taking a matrix exponential; the two agree to $3\times10^{-15}$. The engine uses the closed
form~\eqref{eq:kernel} and keeps the matrix exponential only as a background check.

% ===========================================================================
\section{Tips: what the readout says}\label{sec:tips}
% ===========================================================================

A tip's vector is its emission probability, written on slots:
\begin{align}
\text{tape read as }v_{cz}: &\quad L_{cz}(d)=(1-\pi_{cz})\,\mathbf 1\big[d=|v_{cz}|\big],\qquad L_{cz}(\Sil)=0;\\
\text{tape missing}: &\quad L_{cz}(d)=\pi_{cz}+(1-\pi_{cz})\,\mathbf 1[d=0],\qquad L_{cz}(\Sil)=1 .
\end{align}
A read tape was captured (probability $1-\pi_{cz}$) and is exactly what was read. A missing tape was
lost technically (probability $\pi_{cz}$, whatever it held), or captured but blank (probability
$1-\pi_{cz}$, possible only at depth 0), or silenced (always missing). The blank-tape term is what
makes missingness informative: a missing tape whose relatives show the first slot was already written
cannot have been censored.

% ===========================================================================
\section{Coordinates for the node times}\label{sec:coords}
% ===========================================================================

Node times must be ordered: every node after its parent and before harvest. An optimiser or sampler
that steps freely can break that order. Instead of policing it, the engine uses coordinates in which
every value is a legal tree.

\paragraph{The ratio transform.} Each internal node gets a real number $x_v$, turned into a fraction
by $\sig(x)=1/(1+e^{-x})$. Working in heights $h=1-t$ (founder at $h=1$, cells at $h=0$):
\begin{equation}
h_r=\sig(x_r),\qquad h_v=\sig(x_v)\,h_{\pa(v)}\ \ \text{for every other internal node},\qquad t_v=1-h_v .
\label{eq:ratio}
\end{equation}
Each node sits a fraction $\sig(x_v)$ of the way from the present up to its parent. Since
$0<\sig<1$, every node lands strictly between its parent and the present, for any $x$. For example,
$x_r=0$, $x_u=1$, $x_w=-1$ give $h_r=0.5$, $h_u=0.731\times0.5=0.366$, $h_w=0.269\times0.5=0.134$. Changing
$x_r$ rescales the whole tree; each lower coordinate places a node within the room its parent leaves.

\paragraph{The Jacobian.} How much time one unit of $x$ covers depends on where you are: near
$\sig\approx0.5$ a unit of $x$ moves a node a lot, in the tails very little. The \emph{Jacobian} is that
stretch factor; for the whole tree,
\begin{equation}
\Big|\frac{\partial h}{\partial x}\Big|=\prod_{v\ \text{internal}}h_{\pa(v)}\,\sig(x_v)\big(1-\sig(x_v)\big),
\qquad h_{\pa(r)}=1 .
\label{eq:jac}
\end{equation}
A density over times becomes a density over $x$ by multiplying by this factor. It matters only when
sampling (Section~\ref{sec:objective}). Both the always-legal property and~\eqref{eq:jac} were checked
on 200 random trees.

% ===========================================================================
\section{The tree prior}\label{sec:prior}
% ===========================================================================

The tree prior says how plausible a set of division times is if cells divide at rate $b$, die at rate
$\delta$, and a fraction $\rho$ were sampled at harvest. For a constant-rate birth--death process with
the founder at height 1,
\begin{equation}
p(\mathcal T\mid b,\delta,\rho)\ \propto\ p_1(1)\prod_{v\ \text{internal}}b\,p_1(h_v),\qquad
p_1(h)=\frac{\rho\,(b-\delta)^2e^{-(b-\delta)h}}{\big[\rho b+\big(b(1-\rho)-\delta\big)e^{-(b-\delta)h}\big]^2},
\label{eq:prior}
\end{equation}
conditioned on the clone having its observed $n$ cells. Here $p_1(h)$ is the probability that a
lineage alive at height $h$ leaves exactly one sampled descendant; each internal node contributes ``a
division happened here, and both resulting lineages left sampled descendants.'' The prior does two
jobs. It keeps the ancestors of identical-looking cells from collapsing onto one instant (the
likelihood alone would push their branches to zero length); and it carries information about
calendar time that the edits alone do not. (The normalising constant for conditioning on $n$ is
checked against simulated trees in validation check V8.)

% ===========================================================================
\section{The objective}\label{sec:objective}
% ===========================================================================

By Bayes' rule,
$P(\text{tree},\theta\mid\text{reads})\propto P(\text{reads}\mid\text{tree},\theta)\,
p(\text{tree}\mid b,\delta,\rho)\,p(\theta)$. Taking logs gives the objective the engine
evaluates:
\begin{equation}
J(x,\theta)=\underbrace{\sum_{z}\ell_z\big(t(x),\theta\big)}_{\text{log-likelihood}}
+\underbrace{\log p\big(t(x)\mid b,\delta,\rho\big)}_{\text{log tree prior}}
+\underbrace{\log p(\theta)}_{\text{log parameter priors}}
\ \Big[\ +\log\Big|\frac{\partial h}{\partial x}\Big|\ \text{when sampling}\Big].
\label{eq:J}
\end{equation}
\begin{itemize}
\item \textbf{Point estimate} (maximum a posteriori, i.e.\ maximum likelihood with the priors as a
  penalty): maximise $J$ \emph{without} the Jacobian. Re-expressing where you look does not move the
  highest point of a function, so the best $x$ gives the best times.
\item \textbf{Sampling} the posterior by moving in $x$ (e.g.\ Hamiltonian Monte Carlo): include the
  Jacobian, so the sampler spends time in proportion to probability in \emph{time}, not in $x$.
\end{itemize}
Every step from $x$ to $J$ --- ratio transform, branch totals, kernels, pruning, priors --- is a
smooth formula, so automatic differentiation gives the exact gradient of $J$ with respect to every
node coordinate and parameter.

% ===========================================================================
\section{Per-tape output and support values}\label{sec:support}
% ===========================================================================

Because the engine returns one log-likelihood per tape, any re-weighting
$\ell_{\mathbf w}=\sum_zw_z\ell_z$ is free. A tape bootstrap uses $\mathbf w$ = counts of a
with-replacement draw of $k$ tapes; a jackknife leaves one tape out; a weighted likelihood bootstrap
uses $\mathbf w=k\times\mathrm{Dirichlet}(1,\dots,1)$. The tape is the right resampling unit because
tapes, not slots, are what the model treats as independent. A tape whose $\ell_z$ is persistently out
of line with the others flags a problem tape.

% ===========================================================================
\section{What the engine does not do (yet)}\label{sec:excl}
% ===========================================================================
Read or sequencing errors (would need a larger state representation); cross-tape coupling (shared
cell-level factors; a silenced integration removing its symbol from other tapes); lineage-varying
editing rate or composition (the hooks exist, fixed at ``nothing varies''); partial silencing;
editing coupled to division; editing before the clone founder; pooling across clones; the different
behaviour of slot 6 measured on Park. A \emph{posterior mode} --- posteriors of ancestral states, of
which edits were written on which branch, and of why each missing entry is missing --- is specified
in the design record and deferred to a later version.

% ===========================================================================
\section{Where the pieces come from}\label{sec:prov}
% ===========================================================================

\begin{center}\small
\begin{tabular}{@{}p{0.38\textwidth}p{0.27\textwidth}p{0.29\textwidth}@{}}
\toprule
component & taken from / similar to & new here \\
\midrule
tape as the unit; prefix states; kernel (Eqs 5--7); ancestral sets (Eqs 8--9); pruning (Eq 12) & SciPhy & --- \\
time tree with stem, fixed experiment length, birth--death-sampling prior, per-tape clocks & SciPhy / BEAST & --- \\
absorbing silenced state; editing and readout as two layers & LAML (its kernel equals ours with $N=1$) & per tape, calendar time, timing shape, $\Hbar$ \\
general per-tip emission & LAML-Pro & per (cell, tape) dropout; depth-0 censoring \\
regularising branch lengths & ConvexML's lesson & via the tree prior \\
ratio coordinates & Kishino--Thorne--Bruno; BEAST & --- \\
depth-slot representation (anchored / floating) & extends SciPhy Eqs 8--9 & new (checked exact) \\
time-varying rate and composition via segments & our simulator & new for these methods \\
per-tape log-likelihood array as output & tape bootstraps (Park, Typewriter) & new as an engine output \\
\bottomrule
\end{tabular}
\end{center}

% ===========================================================================
\section{How the code is checked}\label{sec:val}
% ===========================================================================

Two implementations exist: a plain NumPy \emph{reference twin}, written to be read, and the JAX engine,
written to be fast and differentiable. A failed check is a bug, not a finding.
\begin{description}\itemsep2pt
\item[V1] The twin against exhaustive enumeration of every hidden history on small trees.
\item[V2] On the same small settings, the probability of every possible set of reads sums to 1
  (Section~\ref{sec:normalise}).
\item[V3] JAX against the twin on simulated data on library trees.
\item[V4] The closed-form kernel against the matrix exponential.
\item[V5] Automatic gradients against finite differences.
\item[V6] The base model (no dropout or silencing) against SciPhy's own Java likelihood.
\item[V7] The average gradient at the true parameters, over many simulated datasets, is zero.
\item[V8] The tree prior against the library's simulated trees.
\end{description}
V1--V6 are exact up to rounding, each with a stated tolerance. V7--V8 are statistical: the deviation
from zero in standard errors across trees, passing if no deviation exceeds 3.

\subsection*{Results of the first full run (2026-10-06)}
\begin{center}\small
\begin{tabular}{@{}llrrrl@{}}
\toprule
check & compared against & configs & worst & tolerance & verdict\\
\midrule
V1 & full state space / literal enumeration & 340 / 144 & $5.8\times10^{-15}$ & $10^{-10}$ & pass\\
V2 & sum over all possible reads $=1$ & 30 & $1.8\times10^{-14}$ & $10^{-10}$ & pass\\
V3 & JAX vs twin, per tape & 2{,}700 & $1.9\times10^{-15}$ & $10^{-10}$ & pass\\
V4 & closed form vs matrix exponential & 1{,}000 & $1.1\times10^{-14}$ & $10^{-12}$ & pass\\
V5 & automatic vs finite-difference gradients & 360 & $2.5\times10^{-4}$ & $10^{-5}$ & fail (check)\\
V7 & mean score at the truth & 12 & $|z|=1.96$ & $3$ & pass\\
V8 & tree prior vs library trees & 198 & $|z|=3.03$ & $3$ & two trips\\
\bottomrule
\end{tabular}
\end{center}
The likelihood itself is exact: V1--V4 agree to within a few units of the last digit of double
precision. V5 failed on one comparison of 360, a gradient taken at a boundary where silencing is
switched off and the function curves sharply; shrinking the finite-difference step makes the
difference converge to the automatic gradient (error $2.6\times10^{-2}$, $2.5\times10^{-4}$,
$2.5\times10^{-6}$, $2.3\times10^{-8}$ at steps $10^{-4}$ to $10^{-7}$), so the gradient is right and
the check's step was too coarse. V8's two trips are adjacent deciles in a single library cell (20 fixed
trees), unresolved; a follow-up with freshly simulated trees is proposed. The per-tape likelihood costs
$0.04$\,s at 210 cells and $2$\,s at 11{,}081 cells on one CPU core; the gradient grows faster than
linearly at the largest sizes, which a level-by-level implementation would fix.

\end{document}
